(a) The maps reverse inclusions by \exref{I.2.3}. If $\mathfrak a$ is homogeneous and radical, $\mathfrak a\ne S_+$, then either $\mathfrak a=S$, corresponding to the empty set, or $Z(\mathfrak a)\ne\varnothing$ by \exref{I.2.2}. In the latter case $I(Z(\mathfrak a))=\mathfrak a$. Also $Z(I(Y))=Y$ for closed $Y$.

(b) If $Y$ is irreducible and homogeneous $fg$ vanishes on $Y$, then $Y\subseteq Z(f)\cup Z(g)$, so $f$ or $g$ vanishes on $Y$. Thus $S/I(Y)$ has no homogeneous zero divisors, hence no zero divisors by considering highest nonzero homogeneous terms. Conversely, if $Y=Y_1\cup Y_2$ is a union of proper closed subsets, choose homogeneous $f_i\in I(Y_i)\setminus I(Y)$. Then $f_1f_2\in I(Y)$, so $I(Y)$ is not prime.

(c) $I(\mathbb P^n)=(0)$, which is prime.
